% Minimal preamble for Open Logic Project sources (pinned commit 1e960beff9ed7835bf3e3f1335e21af3439cd107).
%
% OLP's own master preamble cannot run on the project host: TeX Live 2025 there lacks
% prooftree.sty, ntheorem.sty, extarrows.sty, mathdots.sty, isomath.sty and subfiles.sty.
% Following the project's own rule ("if the author preamble needs packages the toolchain lacks,
% do not invent macros -- use a minimal working preamble"), this file declares, verbatim in
% spirit and numbering behaviour, exactly the theorem-like environments OLP defines in
% sty/open-logic-book-envs.sty (thm/lem/prop/cor/defn/ex/prob/rem/note/case/axiom/conv),
% plus the shorthand accents OLP uses (`!!a`, `!!b`, `!!{...}`).
%
% Author text is untouched; only the packaging is ours.
\documentclass[10pt,a4paper]{article}

\usepackage{amsmath}
\usepackage{amssymb}
\usepackage{amsthm}
\usepackage{graphicx}
\usepackage{hyperref}
\usepackage{enumitem}
\usepackage[framemethod=default]{mdframed}

% ---- OLP shorthand accents (from sty/open-logic.sty) ----
% NB: \! is taken by LaTeX (negative thin space), so we only define the !!a / !!b families
% and the \olup helper used by them.
\newcommand*{\olup}[1]{\MakeUppercase{#1}}
\def\!!!a#1{\olup{#1}}
\def\!a#1{\olup{#1}}
\def\!b#1{\olup{#1}}

% ---- theorem-like environments, same family numbering as OLP ----
\theoremstyle{plain}
\newtheorem{thm}{Theorem}[section]
\newtheorem{lem}[thm]{Lemma}
\newtheorem{prop}[thm]{Proposition}
\newtheorem{cor}[thm]{Corollary}
\newtheorem{axiom}{Axiom}

\theoremstyle{definition}
\newtheorem{defn}[thm]{Definition}
\newtheorem{ex}[thm]{Example}
\newtheorem{prob}{Problem}[section]
\newtheorem{probd}{Problem}
\newtheorem{conv}[thm]{Convention}

\theoremstyle{remark}
\newtheorem{rem}[thm]{Remark}
\newtheorem{note}[thm]{Note}
\newtheorem{case}[thm]{Case}

% ---- environments OLP defines as wrappers ----
\newenvironment{probtag}[1]{\begin{prob}}{\end{prob}}
\newenvironment{explain}{}{}
\newenvironment{explain*}{}{}
\newenvironment{editorial}{}{}
\newenvironment{digress}{}{}
\newenvironment{digress*}{}{}
\newenvironment{history}{}{}
\newenvironment{reference}{}{}
\newcommand*{\olfileid}[3]{}
\newcommand*{\olsection}[1]{\section*{#1}}
\newcommand*{\olchapter}[3]{}
\newcommand*{\olpart}[2]{}
\newcommand*{\olimport}[1]{}
\newcommand*{\olinclude}[1]{}
\newcommand*{\olpath}[1]{}
\newcommand*{\ollangid}[1]{}
\newcommand*{\problemtag}[1]{}
\newcommand*{\problemauthor}[1]{}
\newcommand*{\ollabel}[1]{\label{#1}}
% \olref takes 0-3 OPTIONAL arguments plus one mandatory label:
%     \olref{lab}  |  \olref[from]{lab}  |  \olref[from][to]{lab}  |  \olref[a][b][c]{lab}
% A single-argument definition silently swallows "[sth]" and emits a broken reference
% (measured on the host: the text came out as "??iso]x"). The chain below consumes the
% brackets and cross-references only the mandatory argument.
\makeatletter
\newcommand{\olref}{\@ifnextchar[{\olref@i}{\olref@final}}
\def\olref@i[#1]{\@ifnextchar[{\olref@ii}{\olref@final}}
\def\olref@ii[#1]{\@ifnextchar[{\olref@iii}{\olref@final}}
\def\olref@iii[#1]{\olref@final}
\def\olref@final#1{\ref{#1}}
\makeatother
\newcommand*{\olnewpage}{\clearpage}

% ===================== OLP compatibility layer (1) =====================

\makeatletter

% ---------------- proof trees (OLP's prooftree is a bussproofs derivative) ----------------
\providecommand{\AxiomC}[1]{}
\providecommand{\UnaryInfC}[1]{}
\providecommand{\BinaryInfC}[1]{}
\providecommand{\TrinaryInfC}[1]{}
\providecommand{\QuaternaryInfC}[1]{}
\providecommand{\QuinaryInfC}[1]{}
\providecommand{\RightLabel}[1]{}
\providecommand{\LeftLabel}[1]{}
\providecommand{\doubleLine}{}
\providecommand{\noLine}{}
\providecommand{\DisplayProof}{}
\providecommand{\fCenter}{}
\newenvironment{prooftree}{}{}
\newenvironment{ltprooftree}{}{}
\newenvironment{oltableau}{}{}

% ---------------- logical notation ----------------
\providecommand{\lif}{\mathbin{\rightarrow}}
\providecommand{\lforall}{\mathord{\forall}}
\providecommand{\lexists}{\mathord{\exists}}
\providecommand{\ltrue}{\top}
\providecommand{\lfalse}{\bot}
\providecommand{\Ltrue}{\top}
\providecommand{\Lfalse}{\bot}
\providecommand{\Sequent}{\mathrel{\Rightarrow}}
\providecommand{\Proves}{\mathrel{\vdash}}
\providecommand{\Entails}{\mathrel{\models}}
\providecommand{\Subst}[3]{#1[#2/#3]}
\providecommand{\sFmla}[2]{#1(#2)}
\providecommand{\pFmla}[2]{#1(#2)}
\providecommand{\formula}{\relax}
\providecommand{\True}{\top}
\providecommand{\False}{\bot}
\providecommand{\Args}[1]{#1}
\providecommand{\Ts}{\ensuremath{\top}}
\providecommand{\Fi}{\ensuremath{\bot}}
\providecommand{\lfalsum}{\bot}

% ---------------- objects, structures, metatheory ----------------
\providecommand{\Obj}{}
\providecommand{\Struct}{\ensuremath{\mathfrak{A}}}
\providecommand{\Domain}{\ensuremath{|\mathfrak{A}|}}
\providecommand{\Assign}{\ensuremath{s}}
\providecommand{\Sat}[2]{\ensuremath{#1 \models #2}}
\providecommand{\mSat}[3]{\ensuremath{#1 \models #2[#3]}}
\providecommand{\mModel}[2]{\ensuremath{#1 \models #2}}
\providecommand{\Lang}{\ensuremath{\mathcal{L}}}
\providecommand{\Th}{\ensuremath{\mathrm{Th}}}
\providecommand{\Log}{\ensuremath{\mathrm{L}}}
\providecommand{\tuple}[1]{\langle #1 \rangle}
\providecommand{\num}[1]{\underline{#1}}
\providecommand{\cfind}[1]{\varphi_{#1}}
\providecommand{\Nat}{\ensuremath{\mathbb{N}}}
\providecommand{\Mult}{\ensuremath{\mathrm{Mult}}}
\providecommand{\Add}{\ensuremath{\mathrm{Add}}}
\providecommand{\Exp}{\ensuremath{\mathrm{Exp}}}
\providecommand{\Zero}{\ensuremath{0}}
\providecommand{\Succ}{\ensuremath{\mathrm{S}}}
\providecommand{\lambd}{\ensuremath{\lambda}}
\providecommand{\bredpar}{\mathrel{\rightarrow_\beta}}
\providecommand{\bcd}[1]{\ensuremath{\ulcorner #1 \urcorner}}
\providecommand{\Set}[1]{\ensuremath{\{#1\}}}
\providecommand{\indcase}[1]{\par\noindent\emph{Case #1}.\quad}
\providecommand{\indfrm}[1]{\par\noindent\emph{Form #1}.\quad}
\providecommand{\iftag}[1]{#1}
\providecommand{\tagitem}[1]{\item[#1]}
\providecommand{\Elim}[1]{}
\providecommand{\Intro}[1]{}
\providecommand{\LeftR}[1]{}
\providecommand{\RightR}[1]{}
\providecommand{\TAss}[1]{}
\providecommand{\TRule}[1]{}

\makeatother
\providecommand{\olpath}{.}

% ===================== OLP compatibility layer (2) =====================
% --- second compatibility block: notation and shorthands used by the frozen OLP units ---
\providecommand{\CutCS}{\mathrm{Cut}}
\providecommand{\LeftSubproofLabel}{\relax}
\providecommand{\RightSubproofLabel}{\relax}
\providecommand{\maeh}{\relax}
\providecommand{\cheight}{\relax}
\providecommand{\ande}{\mathrm{and}_E}
\providecommand{\andi}{\mathrm{and}_I}
\providecommand{\ore}{\mathrm{or}_E}
\providecommand{\ori}{\mathrm{or}_I}
\providecommand{\typeof}{\mathord{:}}
\providecommand{\pValue}[1]{#1}
\providecommand{\Sent}{\mathrm{Sent}}
\providecommand{\cardplus}{\mathbin{+}}
\providecommand{\cardtimes}{\mathbin{\cdot}}
\providecommand{\cardexpo}{\mathbin{\hat{\ }}}
\providecommand{\cardnless}{\not<}
\providecommand{\cardle}{\leq}
\providecommand{\funfromto}[3]{#1\colon #2 \to #3}
\providecommand{\funimage}[1]{\operatorname{im}(#1)}
\providecommand{\fdefined}{\mathord{\downarrow}}
\providecommand{\canonord}{\operatorname{ord}}
\providecommand{\ordtimes}{\mathbin{\cdot}}
\providecommand{\ordeq}[2]{#1 = #2}
\providecommand{\ordneq}[2]{#1 \neq #2}
\providecommand{\ordtype}[1]{\operatorname{ord}(#1)}
\providecommand{\ordsucc}{\mathord{+1}}
\providecommand{\ordunion}{\cup}
\providecommand{\onesphere}{S^1}
\providecommand{\rotationsgroup}{\mathrm{SO}(2)}
\providecommand{\equivrep}[1]{[#1]}
\providecommand{\rlexless}{\prec}
\providecommand{\ZFC}{\mathrm{ZFC}}
\providecommand{\ZF}{\mathrm{ZF}}
\providecommand{\Rat}{\mathbb{Q}}
\providecommand{\Real}{\mathbb{R}}
\providecommand{\PosInt}{\mathbb{Z}^{+}}
\providecommand{\Id}{\mathrm{id}}
\providecommand{\Bin}{\mathrm{Bin}}
\providecommand{\pto}{\rightharpoonup}
\providecommand{\dom}{\operatorname{dom}}
\providecommand{\ran}{\operatorname{ran}}
\providecommand{\len}{\operatorname{len}}
\providecommand{\concat}{\mathbin{{}^\frown}}
\providecommand{\card}[1]{\lvert #1 \rvert}
\providecommand{\Pow}{\mathcal{P}}
\providecommand{\Setabs}[2]{\{\, #1 \mid #2 \,\}}
\providecommand{\comp}{\circ}
\providecommand{\Char}{\chi}
\providecommand{\Prf}{\operatorname{Prf}}
\providecommand{\Prov}{\operatorname{Prov}}
\providecommand{\Proj}{\operatorname{proj}}
\providecommand{\Gn}{\operatorname{gn}}
\providecommand{\gn}{\operatorname{gn}}
\providecommand{\Value}[1]{\operatorname{val}(#1)}
\providecommand{\Deduce}{\mathrm{D}}
\providecommand{\Dual}{\mathrm{d}}
\providecommand{\Ax}{\mathrm{Ax}}
\providecommand{\pSat}[3]{#1 \models #2\,[#3]}
\providecommand{\pAssign}{\sigma}
\providecommand{\bforall}{\forall}
\providecommand{\bexists}{\exists}
\providecommand{\liff}{\mathbin{\leftrightarrow}}
\providecommand{\FV}{\operatorname{FV}}
\providecommand{\UnaryInf}{\ensuremath{\mid}}
\providecommand{\BinaryInf}{\ensuremath{\mid}}
\providecommand{\depth}{\operatorname{depth}}
\providecommand{\oliflabeldef}[3]{#1}
\providecommand{\TMblank}{\sqcup}
\providecommand{\TMstroke}{\mid}
\providecommand{\TMstay}{\mathrm{S}}
\providecommand{\TMright}{\mathrm{R}}
\providecommand{\TMleft}{\mathrm{L}}
\providecommand{\TMendtape}{\triangleright}
\providecommand{\umin}{\mu}
\providecommand{\Sum}{\sum}
\providecommand{\citep}[1]{\cite{#1}}
\providecommand{\citet}[1]{\cite{#1}}
\providecommand{\citeauthor}[1]{\cite{#1}}
\providecommand{\citeyear}[1]{\cite{#1}}

% ===================== OLP compatibility layer (3) =====================
% --- third compatibility block: remaining notation seen in the frozen OLP units ---
\providecommand{\iso}{\cong}
\providecommand{\nszero}{\mathrm{0}}
\providecommand{\nssucc}{\mathrm{S}}
\providecommand{\nsplus}{\mathbin{+}}
\providecommand{\nssless}{<}
\providecommand{\nsless}{<}
\providecommand{\PL}{\mathrm{PL}}
\providecommand{\Nec}{\mathord{\Box}}
\providecommand{\RK}{\mathrm{RK}}
\providecommand{\ST}{\mathrm{ST}}
\providecommand{\SSubst}[3]{#1[#2/#3]}
\providecommand{\subst}[3]{#1[#2/#3]}
\providecommand{\tagfalse}{\mathrm{F}}
\providecommand{\cutr}{\mathrm{Cut}}
\providecommand{\cardneq}{\neq}
\providecommand{\issucc}[1]{\operatorname{succ}(#1)}
\providecommand{\tup}[1]{\langle #1 \rangle}
\providecommand{\Var}{\mathrm{Var}}
\providecommand{\Bew}{\operatorname{Bew}}
\providecommand{\Con}{\operatorname{Con}}
\providecommand{\Cons}{\operatorname{Cons}}
\providecommand{\pDeduce}{\mathrm{D}}
\providecommand{\psat}{\models}
\providecommand{\den}[1]{\llbracket #1 \rrbracket}
\providecommand{\Sem}[2]{\llbracket #1 \rrbracket_{#2}}
\providecommand{\DomainOf}[1]{\lvert #1 \rvert}
\providecommand{\Form}{\mathrm{Form}}
\providecommand{\Terms}{\mathrm{Term}}
\providecommand{\Atom}{\mathrm{Atom}}
\providecommand{\FVars}[1]{\operatorname{FV}(#1)}
\providecommand{\fv}[1]{\operatorname{fv}(#1)}
\providecommand{\tuplep}[1]{\langle #1 \rangle}
\providecommand{\restrict}[2]{#1 \upharpoonright #2}
\providecommand{\sto}{\rightarrow}
\providecommand{\FmlaSet}[1]{\mathrm{Fmla}_{#1}}
\providecommand{\Thm}{\mathrm{Thm}}
\providecommand{\Eval}{\operatorname{val}}
\providecommand{\blank}{\sqcup}
\providecommand{\blankTape}{\sqcup}
\providecommand{\Consis}{\operatorname{Con}}
\providecommand{\drs}{\mathrm{d}}
\providecommand{\VRule}{\mid}
\providecommand{\Bet}{\mathrm{B}}
\providecommand{\Seg}[1]{\mathrm{Seg}_{#1}}
\providecommand{\Enums}{\mathrm{En}}
\providecommand{\Z}{\mathbb{Z}}
\providecommand{\N}{\mathbb{N}}
\providecommand{\abs}[1]{\lvert #1 \rvert}
\providecommand{\floor}[1]{\lfloor #1 \rfloor}
\providecommand{\ceil}[1]{\lceil #1 \rceil}

% ===================== OLP compatibility layer (4) =====================
% --- fourth compatibility block: harvested from the per-candidate compile logs ---
\providecommand{\beredpar}{\mathrel{\rightarrow_\beta}}
\makeatletter
\def\bered{\@ifnextchar[{\bered@i}{\mathrel{\rightarrow_\beta}}}
\def\bered@i[#1]{\mathrel{\rightarrow_\beta}}
\def\bred{\@ifnextchar[{\bred@i}{\mathrel{\rightarrow_\beta}}}
\def\bred@i[#1]{\mathrel{\rightarrow_\beta}}
\makeatother
\providecommand{\beredone}{\mathrel{=_\beta}}
\providecommand{\bredone}{\mathrel{=_\beta}}
\providecommand{\xredone}{\mathrel{=_\beta}}
\providecommand{\aconvone}{\mathrel{=_\alpha}}
\providecommand{\cardless}{<}
\providecommand{\tagrefs}[1]{#1}
\providecommand{\ordplus}{\mathbin{+}}
\providecommand{\funrestrictionto}[2]{#1 \upharpoonright #2}
\providecommand{\fundefined}{\mathord{\uparrow}}
\providecommand{\Trm}{\mathrm{Trm}}
\providecommand{\OProv}{\operatorname{Prov}}
\providecommand{\OPrf}{\operatorname{Prf}}
\providecommand{\aeq}{\mathrel{=}}
\providecommand{\mClass}[1]{[#1]}
\providecommand{\equivclass}[1]{[#1]}
\providecommand{\pheight}{\operatorname{ht}}
\providecommand{\phy}{\operatorname{ht}}
\providecommand{\Cut}{\mathrm{Cut}}
\providecommand{\QuantRank}{\operatorname{qr}}
\providecommand{\elemequiv}{\mathrel{\equiv}}
\providecommand{\PIso}{\mathrm{PIso}}
\providecommand{\Theory}{\mathrm{Th}}
\providecommand{\supstrict}{\subsetneq}
\providecommand{\trcl}{\operatorname{trcl}}
\providecommand{\cref}[1]{\ref{#1}}
\providecommand{\Discharge}{\mathrm{Dis}}
\providecommand{\Weakening}{\mathrm{Wk}}
\providecommand{\ZFminus}{\mathrm{ZF}^{-}}
\makeatletter
\renewcommand{\Dual}{\@ifnextchar[{\Dual@i}{\ifmmode\mathrm{d}\else$\mathrm{d}$\fi}}
\def\Dual@i[#1]{\ifmmode\mathrm{d}[#1]\else$\mathrm{d}[#1]$\fi}
\makeatother
\providecommand{\RightR}{\mathrm{R}}
\providecommand{\LeftR}{\mathrm{L}}
\providecommand{\Intro}{\mathrm{I}}
\providecommand{\Axiom}{\mathrm{Ax}}
\providecommand{\formula}[1]{#1}
\providecommand{\stared}{\ast}
\providecommand{\oset}[2]{\{\, #1 \mid #2 \,\}}
\providecommand{\setof}[2]{\{\, #1 \mid #2 \,\}}
\providecommand{\pow}[1]{\mathcal{P}(#1)}
\providecommand{\ordpair}[2]{\langle #1, #2 \rangle}
\providecommand{\ordtriple}[3]{\langle #1, #2, #3 \rangle}
\providecommand{\union}{\cup}
\providecommand{\inter}{\cap}
\providecommand{\seq}[1]{\langle #1 \rangle}
\providecommand{\SuchThat}{\mid}
\providecommand{\where}{\mid}
\providecommand{\TVar}{\mathrm{TVar}}
\providecommand{\Bool}{\mathrm{Bool}}
\providecommand{\Val}{\operatorname{Val}}
\providecommand{\TL}[1]{\mathrm{TL}}
\providecommand{\xstared}{\ast}

% ===================== OLP compatibility layer (5) =====================
% --- fifth compatibility block: rescues the quarantined OLP candidates ---
\providecommand{\ident}{\operatorname{id}}
\providecommand{\setrank}{\operatorname{rank}}
\providecommand{\cardfont}{\mathfrak{c}}
\providecommand{\Complement}{\overline}
\providecommand{\QR}{\mathrm{QR}}
\providecommand{\becd}{\ensuremath{\ulcorner}}
\providecommand{\disjointsum}{\mathbin{\oplus}}
\providecommand{\cardeq}[2]{#1 = #2}
\providecommand{\xred}{\mathrel{\rightarrow}}
% --- sixth compatibility block: names reported by real compile logs ---
\providecommand{\ORefut}{\operatorname{Refut}}
\providecommand{\ORProv}{\operatorname{Prov}}
\providecommand{\cutrank}{\operatorname{cr}}
\providecommand{\Refut}{\operatorname{Refut}}
\providecommand{\DischargeRule}{\mathrm{Dis}}
\providecommand{\Contraction}{\mathrm{C}}
\providecommand{\bottomAlignProof}{\relax}
\providecommand{\FalseCl}{\bot}
\providecommand{\red}{\mathrel{\rightarrow}}
\providecommand{\varAssign}{\sigma}
\providecommand{\defiff}{\mathrel{\leftrightarrow}}
\providecommand{\FalseInt}{\bot}
\providecommand{\fn}{\operatorname{fn}}
\providecommand{\dotted}{\cdot}

% --- tag-macro semantic fix (2026-10-13, final): \iftag CHOOSES a branch and the OLP proof body is
% always the SECOND group (\iftag{tag}{body}, {tag}{body}{c}, {tag}{body}{c}{d}); a definition that
% prints the tag literal drops the proof body. Bodies contain blank lines, hence \long.
\makeatletter
\providecommand{\iftag}{}
\long\def\iftag#1#2{\@ifnextchar\bgroup{\iftag@three{#2}}{\iftag@stop{#2}}}
\long\def\iftag@three#1#2{\iftag@stop{#1}\@ifnextchar\bgroup{\iftag@eat}{\relax}}
\long\def\iftag@eat#1{}
\long\def\iftag@stop#1{#1}
\long\def\oliflabeldef#1#2{\@ifnextchar\bgroup{\@gobble}{\iftag@stop{#2}}}
\providecommand{\tagrefs}[1]{}
\providecommand{\settag}[1]{}
\makeatother
\begin{document}

% ==== OLP file: content/normal-modal-logic/axioms-systems/axioms-systems.tex ====
% Part: normal-modal-logic
% Chapter: axioms-systems



\olchapter{nml}{prf}{Axiomatic \usetoken{P}{derivation}}

\olimport{introduction}
\olimport{normal-logics}
\olimport{logics-proofs}
\olimport{proofs-in-K}
\olimport{derived-rules}
\olimport{more-proofs-in-K}
\olimport{duals}
\olimport{proofs-modal-systems}
\olimport{soundness}
\olimport{systems-distinct}
\olimport{provability-from-set}
\olimport{provability-properties}
\olimport{consistency}

\OLEndChapterHook



% ==== OLP file: content/normal-modal-logic/axioms-systems/consistency.tex ====
% Part: normal-modal-logic
% Chapter: axioms-systems
% Section: consistency



\olfileid{nml}{prf}{con}

\olsection{Consistency}

Consistency is an important property of sets of \olup{formula}s. A set of
\olup{formula}s is inconsistent if a contradiction, such as~$\lfalse$, is
\olup{derivable} from it; and otherwise consistent. If a set is
inconsistent, its \olup{formula}s cannot all be true in a model at a
world. For the completeness theorem we prove the converse: every
consistent set is true at a world in a model, namely in the ``canonical
model.''

\begin{defn}
  A set $\Gamma$ is \emph{consistent} relatively to a system~$\Sigma$
  or, as we will say, $\Sigma$-consistent, if and only if $\Gamma
  \Proves/[\Sigma] \lfalse$.
\end{defn}

So for instance, the set $\{ \Box(p \lif q), \Box p, \lnot\Box q \}$ is
consistent relatively to propositional logic, but not
\Log{K}-consistent. Similarly, the set $\{ \Diamond p, \Box\Diamond
p \lif q, \lnot q \}$ is not \Log{K5}-consistent.

\begin{prop}\ollabel{prop:consistencyfacts}
  Let $\Gamma$ be a set of \olup{formula}s. Then:
  \begin{enumerate}
  \item $\Gamma$ is $\Sigma$-consistent if and only if there is
    some \olup{formula}~$A$ such that $\Gamma \Proves/[\Sigma]
    A$.
  \item \ollabel{prop:consistencyfacts-b}%
    $\Gamma \Proves[\Sigma] A$ if and only if $\Gamma \cup \{
    \lnot A \}$ is not $\Sigma$-consistent.
  \item \ollabel{prop:consistencyfacts-c}%
    If $\Gamma$ is $\Sigma$-consistent, then for any \olup{formula}
    $A$, either $\Gamma \cup \{ A \}$ is
    $\Sigma$-consistent or $\Gamma \cup \{ \lnot A \}$ is
    $\Sigma$-consistent.
  \end{enumerate}
\end{prop}

\begin{proof}
  These facts follow easily using classical propositional logic. We
  give the argument for \olref{prop:consistencyfacts-c}. Proceed
  contrapositively and suppose neither $\Gamma \cup \{ A \}$ nor
  $\Gamma \cup \{ \lnot A \}$ is $\Sigma$-consistent. Then by
  \olref{prop:consistencyfacts-b}, both $\Gamma, A \Proves[\Sigma]
  \lfalse$ and $\Gamma, \lnot A \Proves[\Sigma] \lfalse$. By the
  deduction theorem $\Gamma \Proves[\Sigma] A \to \lfalse$ and
  $\Gamma \Proves[\Sigma] \lnot A \lif \lfalse$. But $(A \lif
  \lfalse) \lif ((\lnot A \lif \lfalse) \lif \lfalse)$ is a
  tautological instance, hence by
  \olref[prp]{prop:derivabilityfacts}\olref[prp]{prop:derivabilityfacts-ruleT},
  $\Gamma \Proves[\Sigma] \lfalse$.
\end{proof}


% ==== OLP file: content/normal-modal-logic/axioms-systems/derived-rules.tex ====
% Part: normal-modal-logic
% Chapter: axioms-systems
% Section: derived-rules



\olfileid{nml}{prf}{der}

\olsection{Derived Rules}

Finding and writing \olup{derivation}s is obviously difficult,
cumbersome, and repetitive. For instance, very often we want to pass
from $A \lif B$ to $\Box A \lif \Box B$, i.e., apply
rule~\RK. That requires an application of \Nec, then recording the
proper instance of~\Ax{K}, then applying~\MP. Passing from $A \lif
B$ and $B \lif C$ to $A \lif C$ requires recording the (long)
tautological instance
\[
(A \lif B) \lif ((B \lif C) \lif (A \lif C))
\]
and applying \MP{} twice. Often we want to replace a sub-\olup{formula}
by a formula we know to be equivalent, e.g., \iftag{prvDiamond}{$\Diamond
  A$ by $\lnot\Box\lnot A$}{$\Box A$ by $\lnot\Diamond\lnot A$}, or
$\lnot\lnot A$ by $A$. So rather than write out the actual
\olup{derivation}, it is more convenient to simply record why the
intermediate steps are \olup{derivable}. For this purpose, let us collect
some facts about \olup{derivability}.

\begin{prop}
  If $\Log{K} \Proves A_1$, \dots, $\Log{K} \Proves A_n$, and $B$
  follows from $A_1$, \dots, $A_n$ by propositional logic, then
  $\Log{K} \Proves B$.
\end{prop}

\begin{proof}
  If $B$ follows from $A_1$, \dots,~$A_n$ by propositional logic, then
  \[
  A_1 \lif (A_2 \lif \cdots (A_n \lif B)\dots)
  \]
  is a tautological instance. Applying \MP{} $n$ times gives a
  \olup{derivation} of~$B$.
\end{proof}

We will indicate use of this proposition by~\PL.

\begin{prop}
  If $\Log{K} \Proves A_1 \lif (A_2 \lif \cdots (A_{n-1} \lif
  A_n)\dots)$ then $\Log{K} \Proves \Box A_1 \lif (\Box A_2 \lif
  \cdots (\Box A_{n-1} \lif \Box A_n)\dots)$.
\end{prop}

\begin{proof}
  By induction on $n$, just as in the proof of \olref[nor]{prop:rk}.
\end{proof}

We will indicate use of this proposition by~\RK. Let's illustrate how
these results help establishing \olup{derivability} results more easily.

\begin{prop}
  $\Log{K} \Proves (\Box A \land \Box B) \lif \Box (A \land B)$
\end{prop}

\begin{proof}
  \begin{derivation}
    1. & $\Log{K} \Proves A \lif (B \lif (A \land B))$ & \Taut \\
    2. & $\Log{K} \Proves \Box A \lif (\Box B \lif \Box(A \land B)))$ & \RK, 1\\
    3. & $\Log{K} \Proves (\Box A \land \Box B) \lif \Box(A \land B)$ & \PL, 2
  \end{derivation}
\end{proof}

\begin{prop}\ollabel{prop:rewriting}
  If $\Log{K} \Proves A \liff B$ and $\Log{K} \Proves
  \Subst{C}{A}{q}$ then $\Log{K} \Proves \Subst{C}{B}{q}$
\end{prop}

\begin{proof}
  Exercise.
\end{proof}

\begin{prob}
  Prove \olref[nml][prf][der]{prop:rewriting} by proving, by induction
  on the complexity of~$C$, that if $\Log{K} \Proves A \liff B$
  then $\Log{K} \Proves \Subst{C}{A}{q} \liff \Subst{C}{B}{q}$.
\end{prob}

This proposition comes in handy especially when we want to convert
$\Diamond$ into $\Box$ (or vice versa), or remove double negations
inside \olup{formula}.  In what follows, we will mark applications of
\olref{prop:rewriting} by ``$A$ for $B$'' whenever we re-write
\olup{formula}~$C(B)$ for~$C(A)$. In other words, ``$A$ for $B$''
abbreviates:
  \begin{derivation}
    & $\Proves C(A)$ \\
    & $\Proves A \liff B$\\
    & $\Proves C(B)$ & by \olref{prop:rewriting}
  \end{derivation}
For instance:

\begin{prop}
  $\Log{K} \Proves \lnot\Box p \lif \Diamond \lnot p$
\end{prop}

\begin{proof}
  \iftag{prvDiamond}{% Diamond is primitive
  \begin{derivation}
    1. & $\Log{K} \Proves \Diamond \lnot p \liff \lnot\Box\lnot\lnot p$ &
    \Dual\\
    2. & $\Log{K} \Proves \lnot \Box \lnot\lnot p \lif \Diamond \lnot p$ & \PL, 1\\
    3. & $\Log{K} \Proves \lnot\Box p \lif \Diamond \lnot p$ & $p$ for $\lnot\lnot p$\\
  \end{derivation}
  }{% Diamond is defined
  \begin{derivation}
    1. & $\Log{K} \Proves \lnot \Box \lnot\lnot p \lif \Diamond \lnot p$ & \Taut\\
    2. & $\Log{K} \Proves \lnot\Box p \lif \Diamond \lnot p$ & $p$ for $\lnot\lnot p$\\
  \end{derivation}
  The formula on line~$1$ is an instance of $\lnot q \lif \lnot q$,
  substituting $\Box\lnot\lnot p$ for~$q$. $\Diamond\lnot p$ and
  $\lnot\Box\lnot\lnot p$ are identical, since $\Diamond$ is defined
  as~$\lnot\Box\lnot$.}
\end{proof}

In the above \olup{derivation}, the final step ``$p$ for $\lnot\lnot p$''
is short for
  \begin{derivation}
    & $\Log{K} \Proves \lnot \Box \lnot\lnot p \lif \Diamond \lnot
    p$\\
    & $\Log{K} \Proves \lnot\lnot p \liff p$ & \Taut\\
    & $\Log{K} \Proves \lnot\Box p \lif \Diamond \lnot p$ & by \olref{prop:rewriting}
  \end{derivation}
The roles of $C(q)$, $A$, and $B$ in \olref{prop:rewriting} are
played here, respectively, by $\lnot \Box q \lif \Diamond \lnot p$,
$\lnot\lnot p$, and~$p$.

When \olup{formula} contains a sub-\olup{formula} $\lnot\Diamond A$,
we can replace it by $\Box\lnot A$ using \olref{prop:rewriting}, since
$\Log{K} \Proves \lnot\Diamond A \liff \Box\lnot A$. We'll indicate
this and similar replacements simply by ``$\Box\lnot$ for
$\lnot\Diamond$.''

The following proposition justifies that we can establish
\olup{derivability} results schematically. E.g., the previous proposition
does not just establish that $\Log{K} \Proves \lnot\Box p \lif
\Diamond \lnot p$, but $\Log{K} \Proves \lnot\Box A \lif \Diamond
\lnot A$ for arbitrary~$A$.

\begin{prop}
  If $A$ is a substitution instance of $B$ and $\Log{K} \Proves B$,
  then $\Log{K} \Proves A$.
\end{prop}

\begin{proof}
  It is tedious but routine to verify (by induction on the length of
  the \olup{derivation} of~$B$) that applying a substitution to an
  entire \olup{derivation} also results in a correct
  \olup{derivation}. Specifically, substitution instances of tautological
  instances are themselves tautological instances, substitution
  instances of instances of~\iftag{prvDiamond}{\Dual{} and~}{}\Ax{K}
  are themselves instances of~\iftag{prvDiamond}{\Dual{}
    and~}{}\Ax{K}, and applications of \MP{} and~\Nec{} remain correct
  when substituting \olup{formula}s for \olup{propositional variable}s in
  both premise(s) and conclusion.
\end{proof}



% ==== OLP file: content/normal-modal-logic/axioms-systems/duals.tex ====
% Part: normal-modal-logic
% Chapter: axioms-systems
% Section: duals



\olfileid{nml}{prf}{dua}

\olsection{Dual \usetoken{P}{formula}}

\begin{defn}\ollabel{def:duals}
  Each of the \olup{formula}s \Ax{T}, \Ax{B}, \Ax{4}, and
  \Ax{5} has a \emph{dual}, denoted by a subscripted diamond, as
  follows:
  \begin{align}
    \tag{\Ax{T_\Diamond}} p & \lif \Diamond p\\
    \tag{\Ax{B_\Diamond}} \Diamond\Box p & \lif p\\
    \tag{\Ax{4_\Diamond}} \Diamond\Diamond p & \lif \Diamond p\\
    \tag{\Ax{5_\Diamond}} \Diamond\Box p & \lif \Box p
    \end{align}
\end{defn}

Each of the above dual \olup{formula}s is obtained from the corresponding
\olup{formula} by substituting $\lnot p$ for $p$, contraposing, replacing
$\lnot\Box\lnot$ by $\Diamond$, and replacing $\lnot\Diamond\lnot$
by~$\Box$. \Ax{D}, i.e., $\Box A \lif \Diamond A$ is its own dual in
that sense.

\begin{prop}\ollabel{prop:dualsys}
  For each \olup{formula}~$A$ in \olref[nml][prf][dua]{def:duals}:
  $\Log{K}A = \Log{K}A_{\Diamond}$.
\end{prop}
\begin{proof}
  Exercise.
\end{proof}
\begin{prob}
  Prove \olref[nml][prf][dua]{prop:dualsys}.
\end{prob}


% ==== OLP file: content/normal-modal-logic/axioms-systems/introduction.tex ====
% Part: normal-modal-logic
% Chapter: axioms-systems
% Section: introduction



\olfileid{nml}{axs}{int}

\olsection{Introduction}

We have a semantics for the basic modal language in terms of modal
models, and a notion of \olup{formula} being valid---true at all worlds
in all models---or valid with respect to some class of models or
frames---true at all worlds in all models in the class, or based on
the frame. Logic usually connects such semantic characterizations of
validity with a proof-theoretic notion of \olup{derivability}. The aim is
to define a notion of \olup{derivability} in some system such that
\olup{formula} is \olup{derivable} iff it is valid.

The simplest and historically oldest \olup{derivation} systems are
so-called Hilbert-type or axiomatic \olup{derivation} systems.
Hilbert-type \olup{derivation} systems for many modal logics are
relatively easy to construct: they are simple as objects of
metatheoretical study (e.g., to prove soundness and
completeness). However, they are much harder to use to prove
\olup{formula}s in than, say, natural deduction systems.

In Hilbert-type \olup{derivation} systems, \olup{derivation} of \olup{formula}
is a sequence of \olup{formula}s leading from certain axioms, via a
handful of inference rules, to the \olup{formula} in question.  Since we
want the \olup{derivation} system to match the semantics, we have to
guarantee that the set of \olup{derivable} formulas are true in all
models (or true in all models in which all axioms are true). We'll
first isolate some properties of modal logics that are necessary for
this to work: the ``normal'' modal logics. For normal modal logics,
there are only two inference rules that need to be assumed: modus
ponens and necessitation.  As axioms we take all (substitution
instances) of tautologies, and, depending on the modal logic we deal
with, a number of modal axioms. Even if we are just interested in the
class of all models, we must also count all substitution instances
of~$\Ax{K}$\iftag{notprvDiamond}{}{ and~$\Ax{Dual}$} as axioms. This
alone generates the minimal normal modal logic~$\Log K$.

\begin{defn}
The rule of \emph{modus ponens} is the inference schema
\begin{prooftree}
\AxiomC{$A$}
\AxiomC{$A \lif B$}
\RightLabel{\MP}
\BinaryInfC{$B$}
\end{prooftree}
We say \olup{formula}~$B$ \emph{follows from}~\olup{formula}s $A$, $C$
by modus ponens iff $C \ident A \lif B$.
\end{defn}

\begin{defn}
The rule of \emph{necessitation} is the inference schema
\begin{prooftree}
\AxiomC{$A$}
\RightLabel{\Nec}
\UnaryInfC{$\Box A$}
\end{prooftree}
We say the \olup{formula}~$B$ follows from the \olup{formula}s $A$ by
necessitation iff $B \ident \Box A$.
\end{defn}

\begin{defn}
A \emph{\olup{derivation}} from a set of axioms~$\Sigma$ is a sequence of
\olup{formula}s $B_1$, $B_2$, \dots, $B_n$, where each $B_i$ is
either
\begin{enumerate}
\item a substitution instance of a tautology, or
\item a substitution instance of \olup{formula} in~$\Sigma$, or
\item follows from two \olup{formula}s $B_j$, $B_k$ with $j$, $k < i$
  by modus ponens, or
\item follows from \olup{formula}~$B_j$ with $j < i$ by necessitation.
\end{enumerate}
If there is such \olup{derivation} with $B_n \ident A$, we say that
$A$ is \emph{\olup{derivable} from $\Sigma$}, in symbols $\Sigma \Proves A$.
\end{defn}

With this definition, it will turn out that the set of \olup{derivable}
formulas forms a normal modal logic, and that any \olup{derivable}
\olup{formula} is true in every model in which every axiom is true. This
property of \olup{derivation}s is called \emph{soundness}. The converse,
\emph{completeness}, is harder to prove.


% ==== OLP file: content/normal-modal-logic/axioms-systems/logics-proofs.tex ====
% Part: normal-modal-logic
% Chapter: axioms-systems
% Section: logics-proofs



\olfileid{nml}{prf}{prf}

\olsection{\usetoken{P}{derivation} and Modal Systems}

We first define what \olup{derivation} is for normal modal
logics. Roughly, \olup{derivation} is a sequence of \olup{formula}s in
which every \olup{element} is either (a substitution instance of) one of
a number of \emph{axioms}, or follows from previous \olup{element}s by
one of a few inference rules. For normal modal logics, all instances
of tautologies\iftag{prvDiamond}{, \Ax{K}, and~\Dual{}}{ and~\Ax{K}}
count as axioms. This results in the modal system~$\Log{K}$, the
smallest normal modal logic. We may wish to add additional axioms to
obtain other systems, however. The rules are always modus ponens~\MP{}
and necessitation~\Nec.

\begin{defn}
  Given a modal system $\Log{K} A_1 \dots A_n$ and \olup{formula}~$B$
  we say that $B$ is \emph{\olup{derivable}} in $\Log{K} A_1 \dots
  A_n$, written $\Log{K} A_1 \dots A_n \Proves B$, if and only if
  there are \olup{formula}s $C_1$, \dots, $C_k$ such that $C_k = B$
  and each $C_i$ is either a tautological instance, or an instance of
  one of $\Ax{K}$,\iftag{prvDiamond}{ $\Dual$,}{} $A_1$,
  \dots, $A_n$, or it follows from previous \olup{formula}s by means of
  the rules \MP{} or~\Nec.
\end{defn}

The following proposition allows us to show that $B \in \Sigma$
by exhibiting a $\Sigma$-\olup{derivation} of~$B$.

\begin{prop}
  $\Log{K} A_1 \dots A_n = \Setabs{B}{\Log{K} A_1 \dots A_n \Proves B}$.
\end{prop}

\begin{proof}
  We use induction on the length of \olup{derivation}s to show that
  $\Setabs{B}{\Log{K} A_1 \dots A_n \Proves B} \subseteq
  \Log{K} A_1 \dots A_n$.

  If the \olup{derivation} of~$B$ has length~$1$, it contains a single
  \olup{formula}. That \olup{formula} cannot follow from previous formulas
  by \MP{} or \Nec, so must be a tautological instance, an instance of
  \Ax{K},\iftag{prvDiamond}{ $\Dual$,}{} or an instance
  of one of $A_1$, \dots,~$A_n$. But $\Log{K}A_1\dotsA_n$ contains
  these as well, so $B \in \Log{K}A_1 \dots A_n$.

  If the \olup{derivation} of $B$ has length~$> 1$, then $B$ may in
  addition be obtained by \MP{} or \Nec{} from \olup{formula}s not
  occurring as the last line in the \olup{derivation}. If $B$ follows
  from $C$ and $C \lif B$ (by \MP), then $C$ and $C \lif B \in
  \Log{K}A_1 \dots A_n$ by induction hypothesis. But every modal
  logic is closed under modus ponens, so $B \in \Log{K}A_1 \dots
  A_n$. If $B \equiv \Box C$ follows from $C$ by \Nec, then $C
  \in \Log{K}A_1 \dots A_n$ by induction hypothesis. But every
  normal modal logic is closed under $\Nec$, so $B \in
  \Log{K}A_1\dotsA_n$.

  The converse inclusion follows by showing that
  $\Sigma = \Setabs{B}{\Log{K} A_1 \dots A_n \Proves B }$ is a normal
  modal logic containing all the instances of $A_1$,
  \dots, $A_n$, and the observation that $\Log{K} A_1 \dots A_n$
  is, by definition, the smallest such logic.
  \begin{enumerate}
    \item Every tautology~$B$ is a tautological instance, so
      $\Log{K}A_1\dotsA_n \Proves B$, so $\Sigma$ contains all
      tautologies.
    \item If $\Log{K}A_1\dotsA_n \Proves C$ and
      $\Log{K}A_1\dotsA_n \Proves C \lif B$, then
      $\Log{K}A_1\dotsA_n \Proves B$: Combine the \olup{derivation} of
      $C$ with that of $C \lif B$, and add the line~$B$. The last
      line is justified by \MP{}. So $\Sigma$ is closed under modus
      ponens.
    \item If $B$ has \olup{derivation}, then every substitution
      instance of~$B$ also has \olup{derivation}: apply the substitution
      to every \olup{formula} in the \olup{derivation}. (Exercise: prove by
      induction on the length of \olup{derivation}s that the result is
      also a correct \olup{derivation}). So $\Sigma$ is closed under
      uniform substitution. (We have now established that $\Sigma$
      satisfies all conditions of a modal logic.)
    \item We have $\Log{K}A_1\dotsA_n \Proves \Ax{K}$, so $K \in \Sigma$.
    \tagitem{prvDiamond}{We have
      $\Log{K}A_1\dotsA_n \Proves \Dual$, so $\Dual \in \Sigma$.}{}
    \item If $\Log{K}A_1\dotsA_n \Proves C$, the additional
      line~$\Box C$ is justified by~\Nec. Consequently, $\Sigma$ is
      closed under~\Nec. Thus, $\Sigma$ is normal.
  \end{enumerate}
\end{proof}


% ==== OLP file: content/normal-modal-logic/axioms-systems/more-proofs-in-K.tex ====
% Part: normal-modal-logic
% Chapter: axioms-systems
% Section: more-proofs-in-K



\olfileid{nml}{prf}{mpr}

\olsection{More Proofs in \Log{K}}

Let's see some more examples of \olup{derivability} in~\Log{K}, now using
the simplified method introduced in \olref[der]{sec}.

\begin{prop}
  $\Log{K} \Proves \Box (A \lif B) \lif (\Diamond A \lif \Diamond
  B)$
\end{prop}

\begin{proof}
\begin{derivation}
1. & $\Log{K} \Proves (A \lif B) \lif (\lnot B \lif \lnot A)$ & \PL \\
2. &  $\Log{K} \Proves \Box(A \lif B) \lif (\Box\lnot B \lif \Box\lnot A)$ & \RK, 1 \\
3. & $\Log{K} \Proves (\Box\lnot B \lif \Box\lnot A) \lif (\lnot
\Box\lnot A \lif \lnot\Box\lnot B)$ & \Taut\\
4. & $\Log{K} \Proves \Box(A \lif B) \lif (\lnot
\Box\lnot A \lif \lnot\Box\lnot B)$ & \PL, 2, 3 \\
5. &  $\Log{K} \Proves \Box(A \lif B) \lif (\Diamond A \lif \Diamond B)$ & $\Diamond$ for $\lnot\Box\lnot$.
\end{derivation}
\end{proof}
  
\begin{prop}
$\Log{K}\Proves \Box A \lif (\Diamond(A \lif B) \lif
  \Diamond B)$
\end{prop}

\begin{proof}
  \begin{derivation}
  1. & $\Log{K} \Proves A \lif (\lnot B \lif \lnot (A \lif B))$ & \Taut \\
  2. & $\Log{K} \Proves \Box A \lif  (\Box\lnot B \lif \Box\lnot (A\lif B))$ &
  \RK, 1 \\
  3. &  $\Log{K} \Proves \Box A \lif (\lnot \Box\lnot (A\lif B) \lif \lnot
  \Box\lnot B)$ & \PL, 2 \\
  4. & $\Log{K} \Proves \Box A \lif (\Diamond(A \lif B) \lif
  \Diamond B)$ & $\Diamond$ for $\lnot\Box\lnot$.
  \end{derivation}
\end{proof}

\begin{prop}
  $\Log{K} \Proves (\Diamond A \lor \Diamond B) \lif \Diamond(A \lor B)$
\end{prop}

\begin{proof}
  \begin{derivation}
    1. & $\Log{K} \Proves \lnot(A \lor B) \lif \lnot A$ & \Taut \\
    2. & $\Log{K} \Proves \Box\lnot(A \lor B) \lif \Box\lnot A$ & \RK, 1 \\
    3. & $\Log{K} \Proves \lnot\Box\lnot A \lif \lnot\Box\lnot(A \lorB)$ &
    \PL, 2\\
    4. & $\Log{K} \Proves \Diamond A \lif \Diamond(A \lor B)$ &
    $\Diamond$ for $\lnot\Box\lnot$\\
    5. & $\Log{K} \Proves \Diamond B \lif \Diamond(A \lor B)$ & similarly\\
    6. & $\Log{K} \Proves (\Diamond A \lor\Diamond B) \lif \Diamond(A \lor B)$
    & \PL, 4, 5.
  \end{derivation}
\end{proof}

\begin{prop}
  $\Log{K} \Proves \Diamond(A \lorB) \lif (\Diamond A \lor \Diamond B)$
\end{prop}

\begin{proof}
  \begin{derivation}
    1. & $\Log{K} \Proves \lnot A \lif (\lnot B \lif \lnot (A \lor B))$ & \Taut \\
    2. & $\Log{K} \Proves \Box\lnot A \lif
    (\Box\lnot B \lif \Box \lnot (A \lor B))$ & \RK\\
    3. & $\Log{K} \Proves \Box\lnot A \lif (\lnot \Box \lnot (A \lorB)
    \lif \lnot\Box\lnot B)$ & \PL, 2\\
    4. & $\Log{K} \Proves \lnot \Box \lnot(A \lor B) \lif (\Box \lnot A \lif 
    \lnot\Box\lnot B)$ & \PL, 3\\
    5. & $\Log{K} \Proves \lnot \Box \lnot(A\lorB) \lif (\lnot
    \lnot\Box\lnot B \lif \lnot\Box\lnot A)$ & \PL, 4\\
    6. & $\Log{K} \Proves \Diamond(A \lor B) \lif (\lnot
    \Diamond B \lif \Diamond A)$ & $\Diamond$ for $\lnot\Box\lnot$\\
    7. & $\Log{K} \Proves \Diamond(A\lorB) \lif (\Diamond B \lor \Diamond A)$ & \PL, 6. \\
  \end{derivation}
\end{proof}

\begin{prob}
  Show that the following \olup{derivability} claims hold:
  \begin{enumerate}
  \item $\Log{K} \Proves \Diamond \lnot \lfalse \lif (\Box A \lif
    \Diamond A)$;
  \item $\Log{K} \Proves \Box(A \lor B) \lif (\Diamond A \lor \Box
    B)$;
  \item $\Log{K} \Proves (\Diamond A \lif \Box B) \lif \Box(A \lif
    B)$.
  \end{enumerate}
\end{prob}


% ==== OLP file: content/normal-modal-logic/axioms-systems/normal-logics.tex ====
% Part: normal-modal-logic
% Chapter: axioms-systems
% Section: normal-logics



\olfileid{nml}{prf}{nor}

\olsection{Normal Modal Logics}

Not every set of modal \olup{formula}s can easily be characterized as
those \olup{formula}s derivable from a set of axioms. We want modal
logics to be well-behaved. First of all, everything we can \olup{derive}
in classical propositional logic should still be \olup{derivable}, of
course taking into account that the \olup{formula}s may now contain also
\iftag{prvBox}{$\Box$\iftag{prvDiamond}{
    and~}{}}{}\iftag{prvDiamond}{$\Diamond$}{}. To this end, we
require that a modal logic contain all tautological instances and be
closed under modus ponens.

\begin{defn}
  A \emph{modal logic} is a set~$\Sigma$ of modal \olup{formula}s which
  \begin{enumerate}
  \item contains all tautologies, and
  \item is closed under substitution, i.e., if $A \in \Sigma$, and
    $D_1$, \dots, $D_n$ are \olup{formula}s, then
    \[
    \SSubst{A}{\subst{D_1}{p_1}, \dots, \subst{D_n}{p_n}} \in \Sigma,
    \]
    \item is closed under \emph{modus ponens}, i.e., if $A$ and $A
      \lif B \in \Sigma$, then $B \in \Sigma$.
  \end{enumerate}
\end{defn}

In order to use the relational semantics for modal logics, we also
have to require that all \olup{formula}s valid in all modal models are
included. It turns out that this requirement is met as soon as all
instances of \Ax{K}\iftag{prvDiamond}{ and~\Dual{}}{} are
\olup{derivable}, and whenever \olup{formula}~$A$ is \olup{derivable}, so
is~$\Box A$. A modal logic that satisfies these conditions is called
\emph{normal}. (Of course, there are also non-normal modal logics, but
the usual relational models are not adequate for them.)

\begin{defn}
  A modal logic $\Sigma$ is \emph{normal} if it contains
  \begin{align*}
    \tag{\Ax{K}} & \Box(p \lif q) \lif (\Box p \lif \Box q),
    \iftag{prvDiamond}{\\
      \tag{\Dual} & \Diamond p \liff \lnot\Box\lnot p}{}
  \end{align*}
  and is closed under \emph{necessitation}, i.e., if $A \in
  \Sigma$, then $\Box A \in \Sigma$.
\end{defn}

Observe that while tautological implication is ``fine-grained'' enough
to preserve \emph{truth at a world}, the rule \Nec{} only preserves
\emph{truth in a model} (and hence also validity in a frame or in a
class of frames).

\begin{prop}\ollabel{prop:rk}
  Every normal modal logic is closed under rule \RK,
  \begin{prooftree}
    \AxiomC{$A_1 \lif (A_2 \lif \cdots (A_{n-1} \lif A_n)\cdots)$}
    \RightLabel{\RK}
    \UnaryInfC{$\Box A_1 \lif (\Box A_2 \lif \cdots (\Box A_{n-1}
      \lif \Box A_n)\cdots).$}
  \end{prooftree}
\end{prop}

\begin{proof}
  By induction on~$n$: If $n = 1$, then the rule is just \Nec, and
  every normal modal logic is closed under \Nec.

  Now suppose the result holds for $n-1$; we show it holds for~$n$.

  Assume
  \begin{align*}
  & A_1 \lif (A_2 \lif \cdots (A_{n-1} \lif A_n)\cdots) \in \Sigma
  \intertext{By the induction hypothesis, we have}
  & \Box A_1 \lif (\Box A_2 \lif \cdots \Box(A_{n-1} \lif A_n)\cdots)
  \in \Sigma
  \intertext{Since $\Sigma$ is a normal modal logic, it contains all
    instances of~$\Ax{K}$, in particular}
  & \Box(A_{n-1} \lif A_n) \lif (\Box A_{n-1} \lif \Box A_n) \in \Sigma
  \intertext{Using modus ponens and suitable tautological instances we get}
  & \Box A_1 \lif (\Box A_2 \lif \cdots (\Box A_{n-1}
  \lif \Box A_n)\cdots) \in \Sigma. 
  \end{align*}
\end{proof}

\begin{prop}\ollabel{prop:notDiamondBot}
  Every normal modal logic $\Sigma$ contains~$\lnot\Diamond\lfalse$.
\end{prop}

\begin{prob}
  Prove \olref[nml][prf][nor]{prop:notDiamondBot}.
\end{prob}

\begin{prop}
  Let $A_1$, \dots, $A_n$ be \olup{formula}s. Then there is a
  smallest modal logic $\Sigma$ containing all instances of
  $A_1$, \dots, $A_n$.
\end{prop}

\begin{proof}
  Given $A_1$, \dots, $A_n$, define $\Sigma$ as the
  intersection of all normal modal logics containing all instances of
  $A_1$, \dots, $A_n$. The intersection is non-empty as
  $\Frm[L]$, the set of all \olup{formula}s, is such a modal
  logic.
\end{proof}

\begin{defn}
The smallest normal modal logic containing $A_1$, \dots, $A_n$ is
called a \emph{modal system} and denoted by $\Log{K} A_1 \dots
A_n$. The smallest normal modal logic is denoted by~\Log{K}.
\end{defn}


% ==== OLP file: content/normal-modal-logic/axioms-systems/proofs-in-K.tex ====
% Part: normal-modal-logic
% Chapter: axioms-systems
% Section: proofs-in-K



\olfileid{nml}{prf}{prk}

\olsection{Proofs in \Log{K}}

In order to practice proofs in the smallest modal system, we show the
valid \olup{formula}s on the left-hand side of
\olref[syn][sch]{tab:valid-invalidSchemas} can all be given
\Log{K}-proofs.

\begin{prop}
  $\Log{K} \Proves \Box A \lif \Box (B\lif A)$
\end{prop}

\begin{proof}
  \begin{derivation}
    1. & $A \lif (B \lif A)$ & \Taut \\
    2. & $\Box(A \lif (B \lif A))$ & \Nec, 1 \\
    3. & $\Box(A \lif (B \lif A)) \lif
    (\Box A \lif \Box(B \lif A))$ & \Ax{K}\\
    4. & $\Box A \lif \Box (B\lif A)$ & \MP, 2, 3
  \end{derivation}
\end{proof}

\begin{prop}
  $\Log{K} \Proves \Box(A \land B) \lif (\Box A \land \Box B)$
\end{prop}

\begin{proof}
\begin{derivation}
  1. & $(A \land B) \lif A$ & \Taut \\
  2. & $\Box((A \land B) \lif A)$ & \Nec \\
  3. & $\Box((A \land B) \lif A) \lif (\Box(A \land B) \lif \Box A)$
  & \Ax{K} \\
  4. & $\Box(A \land B) \lif \Box A$ & \MP, 2, 3 \\
  5. & $(A \land B) \lif B$ & \Taut \\
  6. & $\Box((A \land B) \lif B)$ & \Nec \\
  7. & $\Box((A \land B) \lif B) \lif (\Box(A \land B) \lif \Box B)$
  & \Ax{K} \\
  8. & $\Box(A \land B) \lif \Box B$ & \MP, 6, 7 \\
  9. & $(\Box(A \land B) \lif \Box A) \lif{}$ \\
  & \qquad $((\Box(A \land B) \lif \Box B) \lif{}$ \\
  & \qquad $(\Box(A \land B) \lif (\Box A \land \Box B)))$ & \Taut\\
  10. &  $(\Box(A \land B) \lif \Box B) \lif{}$ \\
  & \qquad $(\Box(A \land B) \lif (\Box A \land \Box B))$ & \MP, 4, 9\\
  11. & $\Box(A \land B) \lif (\Box A \land \Box B)$ & \MP, 8, 10.
\end{derivation}
Note that the \olup{formula} on line~$9$ is an instance of the tautology
\[
(p \lif q) \lif ((p \lif r) \lif (p \lif (q \land r))).
\]
\end{proof}

\begin{prop}
  $\Log{K}\Proves (\Box A \land \Box B) \lif \Box (A \land B)$
\end{prop}

\begin{proof}
  \begin{derivation}
    1. & $A \lif (B \lif (A \land B))$ & \Taut \\
    2. & $\Box(A \lif (B \lif (A \land B)))$ & \Nec, 1 \\
    3. & $\Box(A \lif (B \lif (A \land B))) \lif (\Box A \lif \Box(B \lif (A \land B)))$ & \Ax{K}\\
    4. & $\Box A \lif \Box(B \lif (A \land B))$ & \MP, 2, 3\\
    5. & $\Box(B \lif (A \land B)) \lif (\Box B \lif \Box(A \land B))$ & \Ax{K} \\
    6. & $(\Box A \lif \Box(B \lif (A \land B))) \lif {}$ \\
    & \qquad $(\Box(B \lif (A \land B)) \lif (\Box B \lif \Box(A \land B))) \lif {}$\\
    & \qquad $(\Box A \lif (\Box B \lif \Box(A \land B))))$ & \Taut\\
    7. & $(\Box(B \lif (A \land B)) \lif (\Box B \lif \Box(A \land B))) \lif {}$\\
    & \qquad $(\Box A \lif (\Box B \lif \Box(A \land B)))$ & \MP, 4, 6\\
    8. & $\Box A \lif (\Box B \lif \Box(A \land B)))$ & \MP, 5, 7\\
    9. & $(\Box A \lif (\Box B \lif \Box(A \land B)))) \lif {}$\\
    & \qquad $((\Box A \land \Box B) \lif \Box(A \land B))$ & \Taut\\
    10. & $(\Box A \land \Box B) \lif \Box(A \land B)$ & \MP, 8, 9
  \end{derivation}
  The \olup{formula}s on lines $6$ and $9$ are instances of the tautologies
  \begin{align*}
    (p \lif q) & \lif ((q \lif r) \lif (p \lif r)) \\
    (p \lif (q \lif r)) & \lif ((p \land q) \lif r)
  \end{align*}
\end{proof}

\begin{prop}
  \iftag{prvBox}{$\Log{K} \Proves \lnot\Box p \lif \Diamond \lnot
    p$}{$\Log{K} \Proves \Box\lnot p \lif \lnot\Diamond p$}
\end{prop}

\begin{proof}
  \iftag{prvBox}{\iftag{prvDiamond}{% Box and Diamond are primitive
  \begin{derivation}
    1. & $\Diamond \lnot p \liff \lnot\Box\lnot\lnot p$ & \Dual\\
    2. & $(\Diamond \lnot p \liff \lnot\Box\lnot\lnot p) \lif {}$ \\
    & \qquad $(\lnot \Box \lnot\lnot p \lif \Diamond \lnot p)$& \Taut\\
    3. & $\lnot \Box \lnot\lnot p \lif \Diamond \lnot p$ & \MP, 1, 2\\
    4. & $\lnot\lnot p \lif p$ & \Taut\\
    5. & $\Box(\lnot\lnot p \lif p)$ & \Nec, 4\\
    6. & $\Box(\lnot\lnot p \lif p) \lif (\Box \lnot\lnot p \lif \Box p)$ & \Ax{K}\\
    7. & $(\Box \lnot\lnot p \lif \Box p)$ & \MP, 5, 6\\
    8. & $(\Box \lnot\lnot p \lif \Box p) \lif (\lnot \Box p \lif \lnot\Box\lnot\lnot p)$ & \Taut\\
    9. & $\lnot \Box p \lif \lnot\Box\lnot\lnot p$ & \MP, 7, 8\\
    10. & $(\lnot \Box p \lif \lnot\Box\lnot\lnot p) \lif {}$ \\
    & \qquad $((\lnot \Box \lnot\lnot p \lif \Diamond \lnot p) \lif (\lnot \Box p \lif \Diamond\lnot p))$ & \Taut\\
    11. & $(\lnot \Box \lnot\lnot p \lif \Diamond \lnot p) \lif (\lnot \Box p \lif \Diamond\lnot p)$ & \MP, 9, 10\\
    12. & $\lnot\Box p \lif \Diamond \lnot p$ & \MP, 3, 11\\
  \end{derivation}
  The \olup{formula}s on lines $8$ and $10$ are instances of the tautologies
  \begin{align*}
    & (p \lif q) \lif (\lnot q \lif \lnot p) \\
    & (p \lif q) \lif ((q \lif r) \lif (p \lif r)).
  \end{align*}
  }{% Diamond is defined
  \begin{derivation}
    1. & $\lnot\lnot p \lif p$ & \Taut\\
    2. & $\Box(\lnot\lnot p \lif p)$ & \Nec, 1\\
    3. & $\Box(\lnot\lnot p \lif p) \lif (\Box \lnot\lnot p \lif \Box p)$ & \Ax{K}\\
    4. & $(\Box \lnot\lnot p \lif \Box p)$ & \MP, 2, 3\\
    5. & $(\Box \lnot\lnot p \lif \Box p) \lif (\lnot \Box p \lif \lnot\Box\lnot\lnot p)$ & \Taut\\
    6. & $\lnot \Box p \lif \lnot\Box\lnot\lnot p$ & \MP, 4, 5\\
    & $\lnot\Box p \lif \Diamond \lnot p$
  \end{derivation}
  The last line is $\lnot \Box p \lif \Diamond\lnot p$, since
  $\Diamond$ is defined as $\lnot\Box\lnot$.  The \olup{formula} on
  line~$5$ is an instance of the tautology
  \[
    (p \lif q) \lif (\lnot q \lif \lnot p).
  \]
  }}{% Box is defined
  \begin{derivation}
    1. & $\Diamond  p \liff \lnot\Box\lnot p$ & \Dual\\
    2. & $(\Diamond p \liff \lnot\Box\lnot p) \lif
    (\Diamond p \lif \lnot \Box \lnot p)$& \Taut\\
    3. & $\Diamond p \lif \lnot \Box \lnot p$ & \MP, 1, 2\\
    4. & $(\Diamond p \lif \lnot \Box \lnot p) \lif (\Box\lnot p \lif \lnot\Diamond p)$ & \Taut\\
    5. & $\Box\lnot p \lif \lnot\Diamond p$ & \MP, 3, 4
  \end{derivation}
  The \olup{formula} on
  line~$4$ is an instance of the tautology
  \begin{align*}
    & (p \lif \lnot q) \lif (q \lif \lnot p).
  \end{align*}
}
\end{proof}

\begin{prob}
  Find \olup{derivation}s in~$\Log{K}$ for the following \olup{formula}s:
  \begin{enumerate}
    \item $\Box \lnot p \lif \Box(p \lif q)$
    \item $(\Box p \lor \Box q) \lif \Box(p \lor q)$
    \item $\Diamond p \lif \Diamond(p \lor q)$
  \end{enumerate}
\end{prob}


% ==== OLP file: content/normal-modal-logic/axioms-systems/proofs-modal-systems.tex ====
% Part: normal-modal-logic
% Chapter: axioms-systems
% Section: proofs-modal-systems



\olfileid{nml}{prf}{prs}

\olsection{Proofs in Modal Systems}

We now come to proofs in systems of modal logic other than~\Log{K}.

\begin{prop}\ollabel{prop:S5facts}
The following provability results obtain:
  \begin{enumerate}
  \item $\Log{KT5} \Proves \Ax{B}$;
  \item $\Log{KT5} \Proves \Ax{4}$;
  \item $\Log{KDB4} \Proves \Ax{T}$;
  \item $\Log{KB4} \Proves \Ax{5}$;
  \item $\Log{KB5} \Proves \Ax{4}$;
  \item \ollabel{prop:S5facts-KT-D}$\Log{KT} \Proves \Ax{D}$.
  \end{enumerate}
\end{prop}

\begin{proof}
  We exhibit proofs for each.
  \begin{enumerate}
    \item $\Log{KT5} \Proves \Ax{B}$:
      \begin{derivation}
        1. & $\Log{KT5} \Proves \Diamond A \lif \Box\Diamond A$ & \Ax{5}\\
        2. & $\Log{KT5} \Proves A \lif \Diamond A$ & $\Ax{T_\Diamond}$\\
        3. & $\Log{KT5} \Proves A \lif \Box\Diamond A$ & \PL, 2, 1.
      \end{derivation}
    \item $\Log{KT5} \Proves \Ax{4}$:
      \begin{derivation}
        1. &$\Log{KT5} \Proves \Diamond\Box A \lif \Box\Diamond\Box A$ & \Ax{5}
        with $\Box A$ for $p$\\
        2. & $\Log{KT5} \Proves \Box A \lif \Diamond\Box A$ & $\Ax{T_\Diamond}$
        with $\Box A$ for $p$\\
        3. & $\Log{KT5} \Proves \Box A \lif \Box\Diamond\Box A$ & \PL, 2, 1\\
        4. & $\Log{KT5} \Proves \Diamond\Box A \lif \Box A$ & $\Ax{5_\Diamond}$\\
        5. & $\Log{KT5} \Proves \Box\Diamond\Box A \lif \Box\Box A$ & \RK{}, 4 \\
        6. & $\Log{KT5} \Proves \Box A \lif \Box\Box A$ & \PL, 3, 5. \\
      \end{derivation}
    \item $\Log{KDB4} \Proves \Ax{T}$:
      \begin{derivation}
        1. & $\Log{KDB4} \Proves \Diamond\Box A \lif A $ & $\Ax{B_\Diamond}$ \\
        2. & $\Log{KDB4} \Proves \Box\Box A \lif \Diamond\Box A$ & $\Ax{D}$
        with $\Box A$ for $p$\\
        3. & $\Log{KDB4} \Proves \Box\Box A \lif A$ & \PL, 1, 2\\
        4. & $\Log{KDB4} \Proves \Box A \lif \Box\Box A$ & \Ax{4} \\
        5. & $\Log{KDB4} \Proves \Box A \lif A$ & \PL, 4, 3. \\
      \end{derivation}
    \item $\Log{KB4} \Proves \Ax{5}$:
      \begin{derivation}
        1. & $\Log{KB4} \Proves \Diamond A \lif \Box \Diamond\Diamond A$ & \Ax{B}
        with $\Diamond A$ for $p$\\
        2. & $\Log{KB4} \Proves \Diamond\Diamond A \lif \Diamond A$ &
        $\Ax{4_\Diamond}$ \\
        3. & $\Log{KB4} \Proves \Box\Diamond\Diamond A \lif \Box \Diamond A$ &
        \RK{}, 2 \\
        4. & $\Log{KB4} \Proves \Diamond A \lif \Box\Diamond A$ & \PL, 1, 3.
      \end{derivation}
    \item $\Log{KB5} \Proves \Ax{4}$:
      \begin{derivation}
        1. & $\Log{KB5} \Proves \Box A \lif \Box\Diamond\Box A$ & \Ax{B} with
        $\Box A$ for $p$ \\
        2. & $\Log{KB5} \Proves  \Diamond\Box A \lif \Box A$ &
        $\Ax{5_\Diamond}$\\
        3. & $\Log{KB5} \Proves  \Box\Diamond\Box A \lif \Box\Box A$ & \RK{}, 2 \\
        4. & $\Log{KB5} \Proves  \Box A \lif \Box\Box A$ & \PL, 1, 3.
      \end{derivation}
    \item $\Log{KT} \Proves \Ax{D}$:
      \begin{derivation}
        1. & $\Log{KT} \Proves \Box A \lif A$ & \Ax{T} \\
        2. & $\Log{KT} \Proves A \lif \Diamond A$ & $\Ax{T_\Diamond}$ \\
        3. & $\Log{KT} \Proves \Box A \lif \Diamond A$ &  \PL, 1, 2
      \end{derivation}
  \end{enumerate}
\end{proof}

\begin{defn}
  Following tradition, we define \Log{S4} to be the system
  \Log{KT4}, and \Log{S5} the system \Log{KTB4}.
\end{defn}

The following proposition shows that the classical system \Log{S5} has
several equivalent axiomatizations. This should not surprise, as the
various combinations of axioms all characterize equivalence relations
(see \olref[frd][es5]{prop:equivalences}).

\begin{prop}\ollabel{prop:S5}
  $\Log{KTB4} = \Log{KT5} = \Log{KDB4} = \Log{KDB5}$.
\end{prop}

\begin{proof}
  Exercise.
\end{proof}

\begin{prob}
  Prove \olref[nml][prf][prs]{prop:S5}.
\end{prob}



% ==== OLP file: content/normal-modal-logic/axioms-systems/provability-from-set.tex ====
% Part: normal-modal-logic
% Chapter: axioms-systems
% Section: provability-from-set



\olfileid{nml}{prf}{prg}

\olsection{\usetoken{S}{derivability} from a Set of \usetoken{P}{formula}}

In \olref[prf][prs]{sec} we defined a notion of provability
of \olup{formula} in a system~$\Sigma$. We now extend this notion to
provability in $\Sigma$ from \olup{formula}s in a set~$\Gamma$.

\begin{defn}\ollabel{defn:Gammaproves}
  !!^a{formula}~$A$ is \olup{derivable} in a system~$\Sigma$ from a set of
  \olup{formula}s $\Gamma$, written $\Gamma \Proves[\Sigma] A$ if and
  only if there are $B_1$, \dots, $B_n \in \Gamma$ such that
  $\Sigma \Proves B_1 \lif (B_2 \lif \cdots (B_n \lif A)
  \cdots)$. 
\end{defn}


% ==== OLP file: content/normal-modal-logic/axioms-systems/provability-properties.tex ====
% Part: normal-modal-logic
% Chapter: axioms-systems
% Section: provability-properties



\olfileid{nml}{prf}{prp}

\olsection{Properties of \usetoken{S}{derivability}}

\begin{prop}\ollabel{prop:derivabilityfacts}
  Let $\Sigma$ be a modal system and $\Gamma$ a set of modal
  \olup{formula}s. The following properties hold:
  \begin{enumerate}
  \item\ollabel{prop:derivabilityfacts-monotonicity} \emph{Monotonicity}: If
    $\Gamma \Proves[\Sigma] A$ and $\Gamma \subseteq \Delta$ then
    $\Delta \Proves[\Sigma] A$;
  \item\ollabel{prop:derivabilityfacts-reflexivity}
    \emph{Reflexivity}: If $A \in \Gamma$ then $\Gamma
    \Proves[\Sigma] A$;
  \item\ollabel{prop:derivabilityfacts-cut} \emph{Cut}: If $\Gamma
    \Proves[\Sigma] A$ and $\Delta \cup \{A\} \Proves[\Sigma] B$
    then $\Gamma \cup \Delta \Proves[\Sigma] B$;
  \item\ollabel{prop:derivabilityfacts-deduction} \emph{Deduction theorem}: $\Gamma \cup \{B\}
    \Proves[\Sigma] A$ if and only if $\Gamma \Proves[\Sigma] B \lif
      A$;
  \item \ollabel{prop:derivabilityfacts-ruleT}% \emph{Rule T}: If
    $\Gamma \Proves[\Sigma] A_1$ and \dots and $\Gamma
    \Proves[\Sigma] A_n$ and $A_1 \to (A_2 \lif \cdots (A_n \lif
    B)\cdots)$ is a tautological instance, then $\Gamma
    \Proves[\Sigma] B$.
  \end{enumerate}
\end{prop}

The proof is an easy exercise. Part
\olref{prop:derivabilityfacts-ruleT} of \olref{prop:derivabilityfacts}
gives us that, for instance, if $\Gamma \Proves[\Sigma] A \lor B$
and $\Gamma \Proves[\Sigma] \lnot A$, then $\Gamma \Proves[\Sigma]
B$. Also, in what follows, we write $\Gamma, A \Proves[\Sigma] B$
instead of $\Gamma \cup \{ A\} \Proves[\Sigma] B$.

\begin{defn}
  A set $\Gamma$ is \emph{deductively closed} relatively to a system
  $\Sigma$ if and only if $\Gamma \Proves[\Sigma] A$ implies
  $A \in \Gamma$.
\end{defn}


% ==== OLP file: content/normal-modal-logic/axioms-systems/soundness.tex ====
% Part: normal-modal-logic
% Chapter: axioms-systems
% Section: soundness



\olfileid{nml}{prf}{snd}

\olsection{Soundness}

!!^a{derivation} system is called sound if everything that can be
\olup{derive}d is valid. When considering modal systems, i.e.,
\olup{derivation}s where in addition to \Ax{K} we can use instances of
some \olup{formula}s $A_1$, \dots,~$A_n$, we want every \olup{derivable}
formula to be true in any model in which $A_1$, \dots, $A_n$ are
true.  

\begin{thm}[Soundness Theorem]\ollabel{thm:soundness}
  If every instance of $A_1$, \dots, $A_n$ is valid in the
  classes of models $\mClass{C}_1$, \dots, $\mClass{C}_n$,
  respectively, then $\Log{K}A_1\dots A_n
  \Proves B$ implies that $B$ is valid in the class of
  models $\mClass{C}_1 \cap \dots \cap \mClass{C}_n$.
\end{thm}

\begin{proof}
  By induction on length of proofs. For brevity, put $\mClass{C} =
  \mClass{C}_1 \cap \dots \cap \mClass{C}_n$.
  \begin{enumerate}
  \item Induction Basis: If $B$ has a proof of length~$1$, then it is
    either a tautological instance, an instance
    of~\Ax{K},\iftag{prvDiamond}{ or of~\Dual{},}{} or an instance of
    one of $A_1$, \dots,~$A_n$. In the first case, $B$ is valid in
    $\mClass{C}$, since tautological instance are valid in \emph{any}
    class of models, by \olref[syn][tau]{prop:valid-taut}. Similarly
    in the second case, by
    \olref[syn][sch]{prop:Kvalid}\iftag{prvDiamond}{ and
      \olref[syn][sch]{prop:Dual-valid}}{}. Finally in the third case,
    since $B$ is valid in $\mClass{C}_i$ and $\mClass{C} \subseteq
    \mClass{C}_i$, we have that $B$ is valid in $\mClass{C}$ as well
    by \olref[syn][val]{prop:subset-class}.
  \item Inductive step: Suppose $B$ has a proof of length $k>1$. If
    $B$ is a tautological instance or an instance of one of $A_1$,
    \dots, $A_n$, we proceed as in the previous step. So suppose $B$ is
    obtained by \MP{} from previous \olup{formula}s $C \lif B$ and
    $C$. Then $C \lif B$ and $C$ have proofs of length $<k$, and
    by inductive hypothesis they are valid in~$\mClass{C}$. By
    \olref[syn][sch]{prop:soundMP}, $B$ is valid in $\mClass{C}$ as
    well. Finally suppose $B$ is obtained by \Nec{} from $C$ (so
    that $B = \BoxC$). By inductive hypothesis, $C$ is valid in
    $\mClass{C}$, and by \olref[syn][val]{prop:Nec-rule} so is~$B$.
  \end{enumerate}
\end{proof}


% ==== OLP file: content/normal-modal-logic/axioms-systems/systems-distinct.tex ====
% Part: normal-modal-logic
% Chapter: axioms-systems
% Section: systems-distinct



\olfileid{nml}{prf}{dis}

\olsection{Showing Systems are Distinct}

In \olref[prs]{sec} we saw how to prove that two systems of modal
logic are in fact the same system. \olref[snd]{thm:soundness} allows
us to show that two modal systems $\Sigma$ and $\Sigma'$ are distinct,
by finding \olup{formula}~$A$ such that $\Sigma' \Proves A$ that fails in
a model of~$\Sigma$.

\begin{prop}
  $\Log{KD} \subsetneq \Log{KT}$
\end{prop}

\begin{proof} This is the syntactic counterpart to the semantic fact
  that all reflexive relations are serial. To show $\Log{KD} \subseteq
  \Log{KT}$ we need to see that $\Log{KD} \Proves B$ implies
  $\Log{KT} \Proves B$, which follows from $\Log{KT} \Proves
  \Log{D}$, as shown in
  \olref[prs]{prop:S5facts}\olref[prs]{prop:S5facts-KT-D}.  To show
  that the inclusion is proper, by Soundness
  (\olref[snd]{thm:soundness}), it suffices to exhibit a model of
  \Log{KD} where \Ax{T}, i.e., $\Box p \lif p$, fails (an
  easy task left as an exercise), for then by Soundness $\Log{KD}
  \Proves/ \Box p \lif p$.
\end{proof}

\begin{prop}
  $\Log{KB} \neq \Log{K4}$. 
\end{prop}

\begin{proof}
  We construct a symmetric model where some instance of \Ax{4} fails;
  since obviously the instance is \olup{derivable} for \Log{K4} but not
  in \Log{KB}, it will follow $\Log{K4} \nsubseteq \Log{KB}$.
  Consider the symmetric model $\mModel{M}$ of \olref{fig:Bnot4}. Since
  the model is symmetric, \Ax{K} and~\Ax{B} are true in~$\mModel{M}$
  (by \olref[syn][sch]{prop:Kvalid} and
  \olref[frd][acc]{thm:soundschemas}, respectively). However,
  $\mSat/{M}{\Box p \lif \Box\Box p}[w_1]$.
\end{proof}

\begin{figure}[htpb]
  \centering
  \begin{tikzpicture}[modal]
    \node[world] (w1) [label=above:\mFalse{p},
    label={below:$\mSat{{}}{\Box p}$\\ $\mSat/{{}}{\Box\Box p}$}] {$w_1$}; 
    \node[world] (w2) [label=above:\mTrue{p},
      label=below:$\mSat/{{}}{\Box p}$, right=of w1]{$w_2$}; 
    \draw[->, bend left] (w1) to (w2); 
    \draw[->, bend left] (w2) to (w1); 
  \end{tikzpicture}
  \caption{A symmetric model falsifying an instance of \Ax{4}.}
  \ollabel{fig:Bnot4}
\end{figure}
 
\begin{thm}\ollabel{thm:KTBnot45}
  $\Log{KTB} \Proves/ \Log{4}$ and $\Log{KTB} \Proves/ \Log{5}$.
\end{thm}

\begin{proof}
  By \olref[frd][acc]{thm:soundschemas} we know that all instances of
  \Ax{T} and \Ax{B} are true in every reflexive symmetric model
  (respectively). So by soundness, it suffices to find a reflexive
  symmetric model containing a world at which some instance of~\Ax{4}
  fails, and similarly for~\Ax{5}.  We use the same model for both
  claims. Consider the symmetric, reflexive model in
  \olref{fig:KTBnot45}.  Then $\mSat/{M}{\Box p \to \Box \Box
    p}[w_1]$, so \Ax{4} fails at $w_1$. Similarly, $\mSat/{M}{\Diamond
    \lnot p \lif \Box \Diamond \lnot p}[w_2]$, so the instance
  of~\Ax{5} with $A = \lnot p$ fails at~$w_2$.
\end{proof}

\begin{figure}[htpb]
  \centering
  \begin{tikzpicture}[modal]
    \node[world] (w1) [label={right:\mTrue{p}},
      label={below: $\mSat{{}}{\Box p}$\\
        $\mSat/{{}}{\Box\Box p}$\\
        $\mSat/{{}}{\Diamond\lnot p}$}] {$w_1$};
    \draw[reflexive above] (w1) to (w1);
    \node[world] (w2) [label={right:\mTrue{p}}, label={below:
        $\mSat{{}}{\Diamond\lnot p}$\\
        $\mSat/{{}}{\Box\Diamond\lnot p}$}, right=of w1] {$w_2$} ;
    \draw[reflexive above] (w2) to (w2);
    \node[world] (w3) [label={right:\mFalse{p}},right=of w2] {$w_3$};
    \draw[reflexive above] (w3) to (w3);
    \draw[->, bend left] (w1) to (w2);  
    \draw[->, bend left] (w2) to (w3);  
    \draw[->, bend left] (w3) to (w2);  
    \draw[->, bend left] (w2) to (w1);  
  \end{tikzpicture}
  \caption{The model for \olref{thm:KTBnot45}.}
  \ollabel{fig:KTBnot45}
\end{figure}

\begin{thm}\ollabel{thm:KD5not4}
  $\Log{KD5} \neq \Log{KT4} = \Log{S4}$.
\end{thm}

\begin{proof}
  By \olref[frd][acc]{thm:soundschemas} we know that all instances of
  \Ax{D} and~\Ax{5} are true in all serial euclidean models.
  So  it suffices to find a serial euclidean model containing a world
  at which some instance of~\Ax{4} fails.  Consider the model of
  \olref{fig:KD5not4}, and notice that $\mSat/{M}{\Box p
  \lif \Box\Box p}[w_1]$. 
\end{proof}

\begin{figure}[t]
  \centering
  \begin{tikzpicture}[modal]
    \node[world] (w2) [label=north west:\mTrue{p}]{$w_2$} ;
    \draw[reflexive left] (w2) to (w2);
    \node[world] (w1) [label=right:\mFalse{p},
      label=below:{$\mSat{{}}{\Box p}, \mSat/{{}}{\Box\Box p}$},
      below right=of w2]{$w_1$};
    \node[world] (w3) [label=north east:\mTrue{p}, above right=of w1]{$w_3$} ;
    \draw[reflexive right] (w3) to (w3);
    \node[world] (w4) [label=right:\mFalse{p},above right=of w2]{$w_4$};
    \draw[reflexive above] (w4) to (w4);
    \draw[->] (w1) to (w2);  
    \draw[->] (w1) to (w3);  
    \draw[->, bend left=15] (w2) to (w3);
    \draw[->, bend left=15] (w3) to (w2);  
    \draw[->, bend left=15] (w2) to (w4);
    \draw[->, bend left=15] (w4) to (w2);  
    \draw[->, bend left=15] (w3) to (w4);
    \draw[->, bend left=15] (w4) to (w3);  
  \end{tikzpicture}
  \caption{The model for \olref{thm:KD5not4}.}
  \ollabel{fig:KD5not4}
\end{figure}

\begin{prob}
  Give an alternative proof of \olref[nml][prf][dis]{thm:KD5not4} using a model
  with $3$ worlds.
\end{prob}

\begin{prob}
  Provide a single reflexive transitive model showing that both
  $\Log{KT4} \Proves/ \Ax{B}$ and $\Log{KT4} \Proves/ \Ax{5}$.
\end{prob}


\end{document}
